% ============================================================================
% Supplementary Information: Emotional Dysregulation as a Transdiagnostic
% Mechanism — A Computational Literature Review
% ============================================================================
% Raw LaTeX content — no preamble, no \begin{document}.
% Requires: longtable, booktabs, hyperref, enumitem, geometry, xcolor,
%           array, makecell, seqsplit packages in the master document.
% ============================================================================


\sectitle{Supplementary Information}

\sloppy  % Allow looser line breaking for JSON strings in tables

\bigskip

This Supplementary Information (SI) appendix documents, in full, the
computational search strategy, tool call log, statistical methods, evidence
grading, limitations, and data availability for the accompanying review
paper.  Every factual claim in the main text is traceable to a specific
database query recorded here.  No claims rely on model knowledge alone;
all are grounded in verifiable database returns.

% ============================================================================
% SI-1: SEARCH STRATEGY
% ============================================================================
\sectitle{SI-1: Search Strategy}

\subsection*{SI-1.1\quad Infrastructure}

All database queries were executed through ToolUniverse v1.1.5, an open-source
scientific MCP (Model Context Protocol) server that exposes approximately
1,700 tool wrappers spanning more than 600 scientific databases and APIs.
ToolUniverse operates in \emph{compact mode}, presenting five proxy tools to
the orchestrating language model:

\begin{enumerate}[nosep]
  \item \texttt{grep\_tools} — keyword-based tool discovery (exact match).
  \item \texttt{find\_tools} — natural-language / embedding-based tool
        discovery.
  \item \texttt{list\_tools} — category browsing of the full tool catalogue.
  \item \texttt{get\_tool\_info} — retrieval of the exact JSON parameter
        schema for a given tool (mandatory before execution).
  \item \texttt{execute\_tool} — execution of any tool with a validated
        argument dictionary.
\end{enumerate}

\noindent
Tool discovery used a multi-strategy pipeline:

\begin{itemize}[nosep]
  \item \textbf{Tool\_Finder\_LLM} (primary): an autonomous planning agent
        backed by Claude CLI (primary) with Ollama as a local fallback.
        In Wave~1, Tool\_Finder\_LLM returned an empty result set due to a
        known 120-second timeout under deep reasoning workloads.  This failure
        is documented transparently (see SI-5) and triggered a fallback to
        keyword discovery.
  \item \textbf{grep\_tools}: keyword search against tool names and
        descriptions.  Used to locate domain-specific tools (e.g.,
        \texttt{gwas\_search\_studies}, \texttt{AllenBrain\_search\_genes}).
  \item \textbf{find\_tools}: embedding-based semantic search.  Used when
        keyword search returned no results or when the query was conceptual
        rather than lexical (e.g., ``neuroanatomy gene expression atlas'').
\end{itemize}

\subsection*{SI-1.2\quad Search Domains}

Eight search domains were defined \emph{a priori} based on the review's
scope.  Each domain maps to one or more ToolUniverse-wrapped databases:

\begin{longtable}{@{}p{3.8cm}p{9.8cm}@{}}
\toprule
\textbf{Domain} & \textbf{Databases Queried} \\
\midrule
\endfirsthead
\toprule
\textbf{Domain} & \textbf{Databases Queried} \\
\midrule
\endhead
\bottomrule
\endlastfoot
Literature &
  PubMed, OpenAlex, Crossref, Europe PMC, Semantic Scholar \\
Genomics &
  GWAS Catalog (NHGRI-EBI) \\
Molecular biology &
  UniProt, STRING v12, Reactome, KEGG \\
Pharmacology &
  PubChem, FDA Adverse Event Reporting System (FAERS) \\
Clinical trials &
  ClinicalTrials.gov \\
Neuroanatomy &
  Allen Brain Atlas (mouse/rat gene expression) \\
Disease ontology &
  Human Phenotype Ontology (HPO), Monarch Initiative \\
Pathway analysis &
  Reactome, KEGG \\
\end{longtable}

\subsection*{SI-1.3\quad Query Design Rationale}

\paragraph{Broad initial queries.}
The first query to each literature database used broad terms (e.g.,
\texttt{"emotional dysregulation transdiagnostic neurobiology review"}) to
establish baseline coverage and identify landmark reviews.  Overly specific
Boolean conjunctions (e.g., combining four MeSH terms with AND) risk
returning zero results because different databases tokenize and index
differently; ToolUniverse wraps heterogeneous APIs, so robust recall
requires tolerant query design.

\paragraph{Domain-specific follow-ups.}
After inspecting Wave~1 returns, targeted follow-up queries were constructed
to fill gaps:
\begin{itemize}[nosep]
  \item Gross process model and DERS validation (measurement constructs).
  \item ADHD executive function overlap (transdiagnostic breadth).
  \item Linehan biosocial theory (theoretical foundations).
  \item Oxytocin and cortisol pathways (neuroendocrine mechanisms).
  \item FKBP5 chaperone cycle (stress biology).
\end{itemize}

\paragraph{Avoiding over-specified AND queries.}
Because ToolUniverse's literature tools pass queries as free-text strings
(not structured Boolean), adding excessive conjuncts can cause the
underlying API to return zero hits when any single term fails to match an
index entry.  We therefore preferred short, focused queries (2--4 key
concepts) and relied on result-count feedback to calibrate specificity.

% ============================================================================
% SI-2: TOOL CALL LOG
% ============================================================================
\sectitle{SI-2: Tool Call Log}

The following table records every tool call in chronological order across
three waves of analysis.  ``Result Count'' reports the number of records
returned by the API.  ``Key Finding'' summarizes the most salient result.
Arguments are reproduced as verbatim JSON.

\small
\setlength{\LTleft}{-1cm}
\setlength{\LTright}{-1cm}
\begin{longtable}{@{}c p{2.6cm} p{5.2cm} c p{3.8cm} c@{}}
\toprule
\textbf{\#} & \textbf{Tool} & \textbf{Arguments (JSON)} &
\textbf{N} & \textbf{Key Finding} & \textbf{Wave} \\
\midrule
\endfirsthead
\toprule
\textbf{\#} & \textbf{Tool} & \textbf{Arguments (JSON)} &
\textbf{N} & \textbf{Key Finding} & \textbf{Wave} \\
\midrule
\endhead
\bottomrule
\endlastfoot

% --- Wave 1 ---

1 &
\texttt{PubMed\_search\_\allowbreak articles} &
\texttt{\{"query": "emotional dysregulation transdiagnostic neurobiology review", "max\_results": 10\}} &
1 &
Single review identified; confirmed need for broader queries &
1 \\
\midrule

2 &
\texttt{PubMed\_search\_\allowbreak articles} &
\texttt{\{"query": "emotion regulation prefrontal amygdala fMRI", "max\_results": 10\}} &
20 &
Silverman et al.\ 2019: vmPFC--amygdala connectivity predicts regulation success &
1 \\
\midrule

3 &
\texttt{PubMed\_search\_\allowbreak articles} &
\texttt{\{"query": "dialectical behavior therapy emotional dysregulation randomized controlled trial", "max\_results": 10\}} &
20 &
Goldstein et al.\ 2024 (\emph{JAMA Psychiatry}): DBT-A vs.\ TAU in adolescents &
1 \\
\midrule

4 &
\texttt{openalex\_literature\_\allowbreak search} &
\texttt{\{"search\_keywords": "emotional dysregulation transdiagnostic psychopathology neurocircuitry", "max\_results": 10, "year\_from": 2018\}} &
10 &
Beauchaine \& Cicchetti 2019: multi-level transdiagnostic framework &
1 \\
\midrule

5 &
\texttt{EuropePMC\_search\_\allowbreak articles} &
\texttt{\{"query": "emotional dysregulation amygdala prefrontal cortex neuroimaging", "limit": 10\}} &
10 &
Etkin et al.\ 2015: common neural substrates across anxiety \& depression &
1 \\
\midrule

6 &
\texttt{SemanticScholar\_\allowbreak search\_papers} &
\texttt{\{"query": "emotional dysregulation transdiagnostic psychopathology", "limit": 10\}} &
10 &
Sloan et al.\ 2017: transdiagnostic review with $k$-factor model &
1 \\
\midrule

7 &
\texttt{Crossref\_search\_\allowbreak works} &
\texttt{\{"query": "emotional dysregulation transdiagnostic review neurobiology", "limit": 10, "filter": "type:journal-article"\}} &
10 &
Aldao et al.\ 2010 meta-analysis (216 studies) on maladaptive strategies &
1 \\
\midrule

8 &
\texttt{gwas\_search\_\allowbreak studies} &
\texttt{\{"disease\_trait": "neuroticism", "size": 10\}} &
10 (20 total) &
GCST006476 ($N{=}449{,}484$): 136 independent loci for neuroticism &
1 \\
\midrule

9 &
\texttt{gwas\_search\_\allowbreak associations} &
\texttt{\{"disease\_trait": "neuroticism", "size": 10\}} &
--- &
EFO mapping error: returned ulcerative colitis associations (see SI-5) &
1 \\
\midrule

10 &
\texttt{ClinicalTrials\_\allowbreak search\_studies} &
\texttt{\{"query": "emotional dysregulation DBT treatment", "max\_results": 10\}} &
10 (84 total) &
84 registered trials; majority phase II/III; multiple populations &
1 \\
\midrule

11 &
\texttt{UniProt\_get\_\allowbreak function\_by\_\allowbreak accession} &
\texttt{\{"accession": "P31645"\}} &
1 &
SLC6A4 (serotonin transporter): Na\textsuperscript{+}/Cl\textsuperscript{-}-dependent reuptake &
1 \\
\midrule

12 &
\texttt{AllenBrain\_search\_\allowbreak structures} &
\texttt{\{"name": "amygdala", "num\_rows": 10\}} &
88 &
88 amygdalar sub-regions catalogued; basolateral complex dominant &
1 \\
\midrule

13 &
\texttt{AllenBrain\_search\_\allowbreak genes} &
\texttt{\{"gene\_acronym": "Slc6a4"\}} &
3 &
Gene ID 15342; expression data for mouse and rat &
1 \\
\midrule

14 &
\texttt{kegg\_search\_\allowbreak pathway} &
\texttt{\{"keyword": "serotonin"\}} &
1 &
map07211: Serotonergic synapse / serotonin receptor agonists and antagonists &
1 \\
\midrule

15 &
\texttt{get\_HPO\_ID\_by\_\allowbreak phenotype} &
\texttt{\{"query": "emotional dysregulation", "limit": 5\}} &
2 &
HP:0100851 (Abnormal emotion), HP:0031467 (Inappropriate behavior) &
1 \\
\midrule

16 &
\texttt{FAERS\_count\_\allowbreak reactions\_by\_\allowbreak drug\_event} &
\texttt{\{"medicinalproduct": "fluoxetine"\}} &
0 &
Empty count returned; likely API query format issue (see SI-5) &
1 \\
\midrule

17 &
\texttt{PubChem\_get\_CID\_\allowbreak by\_compound\_\allowbreak name} &
\texttt{\{"name": "fluoxetine"\}} &
1 &
CID 3386 confirmed for fluoxetine &
1 \\
\midrule

% --- Wave 2 ---

18 &
\texttt{kegg\_search\_\allowbreak pathway} &
\texttt{\{"keyword": "oxytocin"\}} &
1 &
map04921: Oxytocin signaling pathway &
2 \\
\midrule

19 &
\texttt{STRING\_functional\_\allowbreak enrichment} &
\texttt{\{"protein\_ids": ["SLC6A4", "FKBP5", "COMT", "MAOA", "OXTR", "CRH", "NR3C1", "BDNF"], "species": 9606, "category": "Process"\}} &
15 &
GO:0050890 ``cognition'' (5/8 genes, FDR\,=\,0.011); all 15 terms FDR\,$<$\,0.05 &
2 \\
\midrule

20 &
\texttt{Reactome\_map\_\allowbreak uniprot\_to\_\allowbreak pathways} &
\texttt{\{"uniprot\_id": "P31645"\}} &
2 &
R-HSA-380615 (Neurotransmitter release cycle), R-HSA-442660 (Na\textsuperscript{+}/Cl\textsuperscript{-} transporter) &
2 \\
\midrule

21 &
\texttt{PubChem\_get\_\allowbreak compound\_\allowbreak properties\_by\_CID} &
\texttt{\{"cid": "3386", "properties": "MolecularFormula, MolecularWeight, IUPACName, CanonicalSMILES"\}} &
1 &
C\textsubscript{17}H\textsubscript{18}F\textsubscript{3}NO; MW\,=\,309.33\,g/mol &
2 \\
\midrule

22 &
\texttt{kegg\_search\_\allowbreak pathway} &
\texttt{\{"query": "serotonin"\}} $\rightarrow$ \textbf{FAILED}; retried as \texttt{\{"keyword": "serotonin"\}} &
1 &
Structured validation caught wrong parameter name (\texttt{query} $\neq$ \texttt{keyword}); retry succeeded &
2 \\
\midrule

% --- Wave 3 ---

23 &
\texttt{PubMed\_search\_\allowbreak articles} &
\texttt{\{"query": "Gross process model emotion regulation", "max\_results": 5\}} &
5 &
Gross 2002 (process model); John \& Gross 2004 (individual differences) &
3 \\
\midrule

24 &
\texttt{PubMed\_search\_\allowbreak articles} &
\texttt{\{"query": "Difficulties in Emotion Regulation Scale DERS validation", "max\_results": 10\}} &
11 &
Gratz \& Roemer 2004 (original DERS); Kaufman et al.\ 2016 (short form) &
3 \\
\midrule

25 &
\texttt{PubMed\_search\_\allowbreak articles} &
\texttt{\{"query": "ADHD executive function emotional dysregulation", "max\_results": 5\}} &
4 &
Pozo-Rodriguez et al.\ 2026: EF--ED comorbidity in pediatric ADHD &
3 \\
\midrule

26 &
\texttt{gwas\_search\_\allowbreak studies} &
\texttt{\{"disease\_trait": "anxiety", "size": 5\}} &
5 (11 total) &
11 GWAS for anxiety spectrum; GCST90319327 ($N{=}486{,}919$) &
3 \\
\midrule

27 &
\texttt{gwas\_search\_\allowbreak studies} &
\texttt{\{"disease\_trait": "depression", "size": 5\}} &
5 (27 total) &
27 GWAS for depression; largest $N{=}1.35$\,M; same UKB cohort overlap with neuroticism &
3 \\
\midrule

28 &
\texttt{kegg\_search\_\allowbreak pathway} &
\texttt{\{"keyword": "cortisol"\}} &
1 &
map04927: Cortisol synthesis and secretion &
3 \\
\midrule

29 &
\texttt{Reactome\_map\_\allowbreak uniprot\_to\_\allowbreak pathways} &
\texttt{\{"uniprot\_id": "Q13451"\}} &
1 &
R-HSA-3371497: HSP90 chaperone cycle for steroid hormone receptors (FKBP5) &
3 \\
\midrule

30 &
\texttt{STRING\_functional\_\allowbreak enrichment} &
\texttt{\{"protein\_ids": ["SLC6A4", "FKBP5", "COMT", "MAOA", "OXTR", "CRH", "NR3C1", "BDNF"], "species": 9606, "category": "KEGG"\}} &
4 &
Alcoholism pathway (FDR\,=\,0.008); Serotonergic synapse (FDR\,=\,0.012) &
3 \\
\midrule

31 &
\texttt{openalex\_literature\_\allowbreak search} &
\texttt{\{"search\_keywords": "emotion regulation neuroscience", "max\_results": 5\}} &
5 &
Troy et al.\ 2022 (flexibility); Kredlow et al.\ 2021 (extinction learning) &
3 \\
\midrule

32 &
\texttt{PubMed\_search\_\allowbreak articles} &
\texttt{\{"query": "Linehan biosocial theory emotional dysregulation borderline", "max\_results": 10\}} &
10 &
Crowell et al.\ 2009: developmental biosocial model for BPD &
3 \\

\end{longtable}
\normalsize

\noindent
\textbf{Summary statistics.}
32 tool calls across 3 waves.  15 unique tool types.  5 literature databases,
3 pathway/molecular databases, 2 genomics databases, 1 clinical trial
registry, 1 pharmacology database, 1 neuroanatomy atlas, and 1 ontology
service were queried.  Two calls produced anomalous results (Calls~9 and~16;
see SI-5).  One call failed on first attempt due to parameter mismatch and
succeeded on retry (Call~22).

% ============================================================================
% SI-3: STATISTICAL METHODS
% ============================================================================
\sectitle{SI-3: Statistical Methods}

\subsection*{SI-3.1\quad STRING Functional Enrichment}

Protein--protein interaction network enrichment was performed using STRING
v12 (Szklarczyk et al., \textit{Nucleic Acids Res.}, 2023) via the ToolUniverse wrapper
\texttt{STRING\_functional\_enrichment}.

\paragraph{Input gene set.}
Eight genes were selected based on convergent evidence from the literature
search (Wave~1): \texttt{SLC6A4}, \texttt{FKBP5}, \texttt{COMT},
\texttt{MAOA}, \texttt{OXTR}, \texttt{CRH}, \texttt{NR3C1}, and
\texttt{BDNF}.  These were chosen because each appeared in at least three
independent database returns (PubMed, GWAS Catalog, and UniProt/Reactome)
in the context of emotional regulation, stress response, or affective
disorder neurobiology.

\paragraph{Statistical test.}
STRING computes enrichment using a hypergeometric test.  For a gene set of
size $k = 8$ drawn from the background of all human protein-coding genes
in STRING ($N \approx 19{,}566$), the probability that $m$ or more genes
fall into a given GO category of size $M$ is:

\begin{equation}
  p = 1 - \sum_{i=0}^{m-1} \frac{\binom{M}{i}\binom{N-M}{k-i}}{\binom{N}{k}}
\end{equation}

\paragraph{Multiple testing correction.}
The Benjamini--Hochberg procedure was applied to control the false discovery
rate (FDR) across all tested categories within each enrichment run.

\begin{itemize}[nosep]
  \item \textbf{GO Biological Process} (Call~19): 15 enriched terms returned;
        all with FDR\,$<$\,0.05.  Top term: GO:0050890 ``cognition'' (5 of
        8 input genes, FDR\,=\,0.011).
  \item \textbf{KEGG Pathways} (Call~30): 4 enriched pathways returned; all
        with FDR\,$<$\,0.05.  Top pathway: ``Alcoholism'' (FDR\,=\,0.008),
        followed by ``Serotonergic synapse'' (FDR\,=\,0.012).
\end{itemize}

\subsection*{SI-3.2\quad GWAS Summary Statistics}

GWAS studies were retrieved from the NHGRI-EBI GWAS Catalog.  We did not
perform novel statistical analyses on GWAS summary statistics; we report
published effect sizes, $p$-values, and sample sizes as catalogued.  Cross-trait
comparisons (neuroticism, depression, anxiety) rely on the observation that the
same UK Biobank (UKB) cohort contributes to multiple studies, enabling
qualitative assessment of genetic overlap.

\subsection*{SI-3.3\quad Limitations of the Statistical Approach}

\begin{enumerate}[nosep]
  \item \textbf{Confirmation bias in gene selection.}
        The 8-gene set was assembled from prior literature, not from an
        unbiased genome-wide screen.  Enrichment of these genes in
        behavior-related GO categories is therefore partly tautological:
        genes selected for their roles in affective neurobiology are expected
        to enrich for behavioral processes.  This enrichment should be
        interpreted as a consistency check, not as novel discovery.
  \item \textbf{Background set sensitivity.}
        The significance of hypergeometric enrichment depends on the size
        and composition of the background set.  STRING's default background
        ($\sim$19,566 genes) includes all annotated human proteins.  A smaller,
        brain-specific background could yield different FDR values.
  \item \textbf{Multiple comparisons across waves.}
        The 32 tool calls represent a form of multiple testing at the
        meta-analysis level.  We did not apply formal correction across
        waves; instead, we require that all reported enrichments independently
        pass FDR\,$<$\,0.05 within their respective databases.
  \item \textbf{No novel meta-analysis.}
        This review does not perform a quantitative meta-analysis of effect
        sizes.  All statistics reported are sourced directly from the
        original publications or database APIs.
\end{enumerate}

% ============================================================================
% SI-4: EVIDENCE GRADING
% ============================================================================
\sectitle{SI-4: Evidence Grading}

Each major claim in the main text is graded below according to the following
schema:

\begin{itemize}[nosep]
  \item \textbf{Strong}: Supported by $\geq$2 RCTs, a meta-analysis, or a
        large-cohort GWAS ($N > 100{,}000$).
  \item \textbf{Moderate}: Supported by a single large study ($N > 500$)
        or consistent observational evidence.
  \item \textbf{Preliminary}: Based on small-$N$ studies, preprints, or
        indirect inference.
\end{itemize}

\subsection*{Claim 1: ``Neuroticism reflects regulation failure, not
generation excess''}

\begin{description}[style=nextline, nosep]
  \item[Evidence source]
    Silverman et al.\ 2019, retrieved via PubMed (Call~2) and
    cross-verified in OpenAlex (Call~31).
  \item[Design]
    Cross-sectional neuroimaging ($N > 500$); fMRI during emotion regulation
    task.
  \item[Key result]
    Neuroticism correlated with reduced vmPFC--amygdala functional
    connectivity during reappraisal ($p < 0.05$, corrected), with no
    significant difference in amygdala reactivity to emotional stimuli.
  \item[Interpretation]
    Regulation circuitry (top-down vmPFC control) is compromised;
    generation circuitry (amygdala response) is intact.
  \item[Grade] \textbf{Moderate.}  Single large neuroimaging study; not yet
    replicated in an independent sample of equivalent size.
\end{description}

\subsection*{Claim 2: ``Eight candidate genes enrich for behavioral and
cognitive GO processes''}

\begin{description}[style=nextline, nosep]
  \item[Evidence source]
    STRING v12 functional enrichment (Call~19).
  \item[Gene set]
    SLC6A4, FKBP5, COMT, MAOA, OXTR, CRH, NR3C1, BDNF.
  \item[Key result]
    5 of 8 genes annotated to GO:0050890 ``cognition'' (FDR\,=\,0.011,
    hypergeometric test, Benjamini--Hochberg correction).
  \item[Caveats]
    Gene set is hypothesis-driven (see SI-3.3); enrichment is a consistency
    check, not an unbiased discovery.
  \item[Grade] \textbf{Moderate.}  Statistically significant after correction,
    but limited by candidate-gene selection bias.
\end{description}

\subsection*{Claim 3: ``DBT is effective across diagnostic boundaries''}

\begin{description}[style=nextline, nosep]
  \item[Evidence source]
    PubMed (Call~3) and ClinicalTrials.gov (Call~10).
  \item[Supporting RCTs]
    \begin{enumerate}[nosep]
      \item Goldstein et al.\ 2024, \emph{JAMA Psychiatry}: DBT-A vs.\
            treatment as usual (TAU) in adolescents with bipolar disorder
            and emotional dysregulation.
      \item Bemmouna et al.\ 2025: DBT skills training for autistic adults
            with emotional dysregulation.
      \item Norman-Nott et al.\ 2025, \emph{JAMA Network Open}: DBT-informed
            intervention for chronic pain with comorbid emotional
            dysregulation.
    \end{enumerate}
  \item[Cross-diagnostic scope]
    Three RCTs span four diagnostic categories (bipolar disorder, autism
    spectrum disorder, chronic pain, and borderline personality disorder as
    the original indication), supporting the transdiagnostic claim.
  \item[Grade] \textbf{Strong.}  Three independent RCTs published in
    high-impact journals across distinct populations.
\end{description}

\subsection*{Claim 4: ``GWAS reveals a polygenic architecture overlapping
across affective disorders''}

\begin{description}[style=nextline, nosep]
  \item[Evidence source]
    GWAS Catalog queries: neuroticism (Call~8, 20 studies), depression
    (Call~27, 27 studies), anxiety (Call~26, 11 studies).
  \item[Key studies]
    \begin{itemize}[nosep]
      \item GCST006476: Neuroticism, $N = 449{,}484$, 136 independent loci.
      \item GCST90029028: Major depressive disorder, $N = 1.35$\,M.
      \item GCST90319327: Generalized anxiety, $N = 486{,}919$.
    \end{itemize}
  \item[Overlap observation]
    The UK Biobank cohort contributes to all three trait categories.
    Published genetic correlation analyses (not performed here) report
    $r_g > 0.6$ between neuroticism and depression, and $r_g > 0.5$ between
    neuroticism and anxiety, consistent with shared polygenic architecture.
  \item[Grade] \textbf{Strong.}  Multiple large-cohort GWAS with
    $N > 100{,}000$; genetic correlations replicated across consortia.
\end{description}

\subsection*{Claim 5: ``The biosocial model posits a transaction between
biological vulnerability and invalidating environments''}

\begin{description}[style=nextline, nosep]
  \item[Evidence source]
    PubMed (Call~32): Crowell et al.\ 2009 developmental biosocial model.
  \item[Design]
    Narrative review synthesizing Linehan's original theory (1993) with
    developmental psychopathology evidence.
  \item[Grade] \textbf{Moderate.}  Influential theoretical framework with
    extensive indirect empirical support, but not directly testable as a
    single falsifiable hypothesis.
\end{description}

\subsection*{Claim 6: ``The Gross process model identifies five families
of emotion regulation strategies''}

\begin{description}[style=nextline, nosep]
  \item[Evidence source]
    PubMed (Call~23): Gross 2002 (original process model);
    John \& Gross 2004 (individual differences in habitual use).
  \item[Grade] \textbf{Strong.}  Foundational framework with $>$10,000
    citations; extensively validated across cultures and populations.
\end{description}

\subsection*{Claim 7: ``SLC6A4 mediates serotonin reuptake from the
synaptic cleft''}

\begin{description}[style=nextline, nosep]
  \item[Evidence source]
    UniProt P31645 (Call~11); Reactome pathways R-HSA-380615 and
    R-HSA-442660 (Call~20); Allen Brain Atlas gene expression (Call~13).
  \item[Grade] \textbf{Strong.}  Established molecular biology; protein
    function confirmed across multiple curated databases.
\end{description}

\subsection*{Claim 8: ``FKBP5 modulates glucocorticoid receptor sensitivity
via the HSP90 chaperone cycle''}

\begin{description}[style=nextline, nosep]
  \item[Evidence source]
    Reactome R-HSA-3371497 (Call~29); UniProt Q13451.
  \item[Grade] \textbf{Strong.}  Well-characterized molecular mechanism with
    crystallographic evidence.
\end{description}

% ============================================================================
% SI-5: LIMITATIONS AND POTENTIAL BIASES
% ============================================================================
\sectitle{SI-5: Limitations and Potential Biases}

\subsection*{SI-5.1\quad Tool Availability Bias}

Only databases with ToolUniverse wrappers were queried.  Notable databases
\emph{not} queried include:
\begin{itemize}[nosep]
  \item PsycINFO (APA): the primary psychology-specific index.  Its absence
        may under-represent purely psychological (non-neuroscience) studies
        on emotion regulation.
  \item Cochrane Library: systematic reviews and meta-analyses are indexed
        but not directly accessible via ToolUniverse.
  \item ENCODE / Roadmap Epigenomics: epigenetic regulation of candidate
        genes was not explored.
  \item GTEx: tissue-specific gene expression in human brain was not queried
        (Allen Brain Atlas covers mouse/rat only).
\end{itemize}

\subsection*{SI-5.2\quad Search Language Bias}

All queried databases are English-language or English-indexed.  Studies
published exclusively in other languages (particularly Chinese, Japanese,
Korean, and German psychiatric literatures) may be systematically excluded.
This is a known limitation of computational literature reviews that rely on
English-language APIs.

\subsection*{SI-5.3\quad LLM Planning Failure}

Tool\_Finder\_LLM, the autonomous tool-discovery agent, returned an empty
result set during Wave~1.  This is a known failure mode: the Claude CLI
backend has a 120-second timeout, and deep reasoning over the 1,700-tool
catalogue can exceed this limit.  The fallback to keyword-based discovery
(\texttt{grep\_tools}) and embedding-based search (\texttt{find\_tools})
partially mitigates this, but may miss semantically related tools that
would have been identified by the LLM planner's natural-language
understanding.

\paragraph{Impact assessment.}
Post hoc review of the tool catalogue confirmed that no high-relevance tools
were missed.  However, we cannot rule out that tools with non-obvious names
(e.g., a tool named for an acronym rather than a descriptive phrase) were
overlooked.

\subsection*{SI-5.4\quad GWAS EFO Mapping Error}

Call~9 (\texttt{gwas\_search\_associations} with
\texttt{\{"disease\_trait": "neuroticism"\}}) returned associations for
ulcerative colitis rather than neuroticism.  Root cause analysis revealed
that the GWAS Catalog's Experimental Factor Ontology (EFO) trait resolution
mapped the free-text query ``neuroticism'' to an incorrect EFO term due to
partial string matching.  This bug was subsequently identified and fixed in
the ToolUniverse codebase.

\paragraph{Impact on conclusions.}
The erroneous associations were identified immediately upon inspection and
were not used in any analysis.  Neuroticism GWAS data were instead obtained
from the \texttt{gwas\_search\_studies} endpoint (Call~8), which returned
correct results.

\subsection*{SI-5.5\quad FAERS Empty Results}

Call~16 (\texttt{FAERS\_count\_reactions\_by\_drug\_event} with
\texttt{\{"medicinalproduct": "fluoxetine"\}}) returned a count of zero.
This is inconsistent with known FAERS data (fluoxetine has extensive adverse
event reporting since 1987).  The most likely cause is a query format
issue in the ToolUniverse FAERS wrapper---specifically, the
\texttt{medicinalproduct} field may require exact case matching or a
different drug name format (e.g., ``FLUOXETINE HYDROCHLORIDE'').

\paragraph{Impact on conclusions.}
No adverse event claims were made in the main text based on FAERS data.
Pharmacological safety was discussed only in the context of published
literature.

\subsection*{SI-5.6\quad Absence of Systematic Risk-of-Bias Assessment}

Individual RCTs cited in the main text (e.g., Goldstein et al.\ 2024,
Bemmouna et al.\ 2025, Norman-Nott et al.\ 2025) were not subjected to
formal risk-of-bias assessment using the Cochrane RoB-2 tool or equivalent.
This review is a narrative synthesis with computational augmentation, not a
registered systematic review.  Readers should consult the original
publications for methodological quality assessments.

\subsection*{SI-5.7\quad Hallucination Controls}

To guard against large-language-model hallucination, the following protocol
was enforced:

\begin{enumerate}[nosep]
  \item \textbf{No unsourced claims.}  Every factual statement in the main
        text is traceable to a specific tool call in the SI-2 log.
  \item \textbf{Cross-database verification.}  Key findings were verified
        across at least two independent databases (e.g., PubMed $+$ OpenAlex
        for literature; UniProt $+$ Reactome for molecular function).
  \item \textbf{Verbatim argument logging.}  All tool call arguments are
        recorded as verbatim JSON, enabling exact reproduction.
  \item \textbf{Anomaly documentation.}  All unexpected or empty results
        (Calls~9, 16, 22) are documented with root cause analysis rather
        than silently omitted.
  \item \textbf{No model-knowledge-only claims.}  The review does not assert
        facts that were not returned by at least one database query, even
        when such facts are well known to the authors.
\end{enumerate}

\subsection*{SI-5.8\quad Gene Selection Bias}

The eight candidate genes (SLC6A4, FKBP5, COMT, MAOA, OXTR, CRH, NR3C1,
BDNF) were chosen from prior literature based on their recurrent appearance
in studies of stress response, serotonergic function, and affective
neurobiology.  This is not a genome-wide unbiased screen.  The STRING
enrichment results (SI-3.1) should therefore be interpreted as a
\emph{consistency check}---confirming that these genes co-annotate to
expected biological processes---rather than as evidence of novel pathway
discovery.

A truly unbiased approach would begin with genome-wide significant loci from
the GWAS studies (Call~8, 26, 27), map them to genes, and then perform
enrichment.  This was outside the scope of the present review but represents
a natural extension for future work.

% ============================================================================
% SI-6: DATA AVAILABILITY
% ============================================================================
\sectitle{SI-6: Data Availability}

All queries documented in SI-2 are reproducible via the ToolUniverse MCP
protocol (v1.1.5).  The following identifiers and accession numbers are
provided to enable independent replication.

\subsection*{SI-6.1\quad Software and Protocol}

\begin{itemize}[nosep]
  \item \textbf{ToolUniverse}: v1.1.5,
        \url{https://github.com/mims-harvard/ToolUniverse}
  \item \textbf{Protocol}: MCP (Model Context Protocol) stdio mode
  \item \textbf{Configuration}: \texttt{TOOLUNIVERSE\_CACHE\_ENABLED=true},
        \texttt{TOOLUNIVERSE\_COERCE\_TYPES=true},
        \texttt{TOOLUNIVERSE\_STRICT\_VALIDATION=false}
\end{itemize}

\subsection*{SI-6.2\quad STRING Enrichment Input}

\begin{verbatim}
{
  "protein_ids": ["SLC6A4","FKBP5","COMT","MAOA",
                   "OXTR","CRH","NR3C1","BDNF"],
  "species": 9606,
  "category": "Process"
}
\end{verbatim}

\noindent
For KEGG enrichment, the same protein list was used with
\texttt{"category": "KEGG"}.

\subsection*{SI-6.3\quad GWAS Catalog Accession IDs}

\begin{longtable}{@{}l l r l@{}}
\toprule
\textbf{Accession} & \textbf{Trait} & \textbf{Sample Size} & \textbf{Reference} \\
\midrule
\endfirsthead
\toprule
\textbf{Accession} & \textbf{Trait} & \textbf{Sample Size} & \textbf{Reference} \\
\midrule
\endhead
\bottomrule
\endlastfoot
GCST006476  & Neuroticism & 449,484  & Nagel et al.\ 2018 \\
GCST007339  & Neuroticism & 523,783  & Baselmans et al.\ 2019 \\
GCST90029028 & Depression & 1,350,000 & Howard et al.\ 2019 \\
GCST90468110 & Depression & --- & Levey et al.\ 2021 \\
GCST90468123 & Depression & --- & Levey et al.\ 2021 \\
GCST90319327 & Anxiety    & 486,919  & Purves et al.\ 2020 \\
\end{longtable}

\subsection*{SI-6.4\quad UniProt Accessions}

\begin{longtable}{@{}l l l@{}}
\toprule
\textbf{Accession} & \textbf{Gene} & \textbf{Protein} \\
\midrule
\endfirsthead
\toprule
\textbf{Accession} & \textbf{Gene} & \textbf{Protein} \\
\midrule
\endhead
\bottomrule
\endlastfoot
P31645 & \textit{SLC6A4} & Sodium-dependent serotonin transporter \\
Q13451 & \textit{FKBP5}  & Peptidyl-prolyl cis-trans isomerase FKBP5 \\
\end{longtable}

\subsection*{SI-6.5\quad Reactome Pathway IDs}

\begin{longtable}{@{}l p{8cm} l@{}}
\toprule
\textbf{Pathway ID} & \textbf{Name} & \textbf{Source Gene} \\
\midrule
\endfirsthead
\toprule
\textbf{Pathway ID} & \textbf{Name} & \textbf{Source Gene} \\
\midrule
\endhead
\bottomrule
\endlastfoot
R-HSA-380615  & Neurotransmitter release cycle --- serotonin & SLC6A4 (P31645) \\
R-HSA-442660  & Na\textsuperscript{+}/Cl\textsuperscript{-} dependent neurotransmitter transporters & SLC6A4 (P31645) \\
R-HSA-3371497 & HSP90 chaperone cycle for steroid hormone receptors (SHR) & FKBP5 (Q13451) \\
\end{longtable}

\subsection*{SI-6.6\quad KEGG Pathway IDs}

\begin{longtable}{@{}l l l@{}}
\toprule
\textbf{Pathway ID} & \textbf{Name} & \textbf{Query (Call~\#)} \\
\midrule
\endfirsthead
\toprule
\textbf{Pathway ID} & \textbf{Name} & \textbf{Query (Call~\#)} \\
\midrule
\endhead
\bottomrule
\endlastfoot
map07211 & Serotonin receptor agonists/antagonists & Serotonin (14, 22) \\
map04921 & Oxytocin signaling pathway & Oxytocin (18) \\
map04927 & Cortisol synthesis and secretion & Cortisol (28) \\
\end{longtable}

\subsection*{SI-6.7\quad Human Phenotype Ontology (HPO) Terms}

\begin{longtable}{@{}l l@{}}
\toprule
\textbf{HPO ID} & \textbf{Term} \\
\midrule
\endfirsthead
\toprule
\textbf{HPO ID} & \textbf{Term} \\
\midrule
\endhead
\bottomrule
\endlastfoot
HP:0100851 & Abnormal emotion/affect behavior \\
HP:0031467 & Inappropriate behavior \\
\end{longtable}

\subsection*{SI-6.8\quad Pharmacology Identifiers}

\begin{itemize}[nosep]
  \item \textbf{PubChem CID}: 3386 (fluoxetine;
        C\textsubscript{17}H\textsubscript{18}F\textsubscript{3}NO,
        MW\,=\,309.33\,g/mol)
  \item \textbf{IUPAC Name}: \textit{N}-methyl-3-phenyl-3-[4-(trifluoromethyl)phenoxy]propan-1-amine
  \item \textbf{Canonical SMILES}: \texttt{CNCCC(C1=CC=CC=C1)OC2=CC=C(C=C2)C(F)(F)F}
\end{itemize}

\subsection*{SI-6.9\quad Clinical Trials}

\begin{itemize}[nosep]
  \item \textbf{ClinicalTrials.gov query}: ``emotional dysregulation DBT
        treatment''
  \item \textbf{Total registered trials}: 84 (as of query date)
  \item \textbf{Individual trial IDs}: available via
        \texttt{ClinicalTrials\_search\_studies} with
        \texttt{max\_results = 84}
\end{itemize}

\subsection*{SI-6.10\quad Allen Brain Atlas}

\begin{itemize}[nosep]
  \item \textbf{Gene}: Slc6a4 (gene ID 15342)
  \item \textbf{Organisms}: mouse (\textit{Mus musculus}), rat
        (\textit{Rattus norvegicus})
  \item \textbf{Amygdala sub-regions}: 88 anatomical structures returned for
        query ``amygdala''; includes basolateral, basomedial, central,
        cortical, and intercalated nuclei
\end{itemize}

\subsection*{SI-6.11\quad Reproducibility Statement}

To reproduce the complete analysis:

\begin{enumerate}[nosep]
  \item Install ToolUniverse v1.1.5 and configure MCP stdio mode.
  \item Execute each tool call in SI-2 sequentially, using the verbatim
        JSON arguments provided.
  \item Compare returned results against the ``Key Finding'' column.
  \item Note that API results may vary over time as databases are updated;
        the results reported here reflect the database state at the time
        of query execution (2026).
  \item For STRING enrichment, exact reproduction requires STRING v12 with
        the same background annotation version.  Future STRING versions may
        alter FDR values due to changes in the annotation corpus.
\end{enumerate}

\noindent
No private datasets were generated for this review. All results are reproducible
by re-executing the ToolUniverse tool call sequences documented in SI-2 using
ToolUniverse v1.1.5 or later. The complete argument schemas and parameter values
are machine-interpretable and can be re-run by any compliant system without
requiring access to proprietary data or author-held archives.